\comoleerla{dos paneles (red corregida y con la longitud registrada de Chocontá a Tunja); cada punto es un $(w_1, w_2)$ de la malla con $w_3 = 1 - w_1 - w_2$; círculo azul, gana Bogotá; triángulo naranja, gana Santander; la línea es la frontera exacta.}
{Bogotá gana en 17 de 66 puntos en la red corregida y en 19 con la longitud registrada; la frontera corta $w_3=0$ en $w_1 = 0.4422$ y $0.4095$.}
{Medido: Bogotá es más corta ($\Delta D = -2.1406$) pero más cara en peaje ($\Delta P = 1.6971$) y en riesgo ($\Delta R = 2.8788$); la frontera es el lugar donde esos tres términos se compensan.}
